\chapter{Métodos e procedimentos}
\label{chap:metodologia}

O presente trabalho baseia-se na estratégia empírica que retoma a função de reação fiscal desenvolvida por \citeonline{abubakarDebtReliefFiscal2025} para 37 países localizados na África Subsaariana durante o período 2000-2019, adaptando-a para o contexto dos estados brasileiros. A adaptação preserva a lógica de identificação via interação entre dívida defasada e variável institucional, com três modificações principais.

Em relação às variáveis fiscais, são expressas em proporção da Receita Corrente Líquida (RCL), e não do PIB, por ser essa a referência adotada na LRF e nos contratos de refinanciamento. A variável institucional, por sua vez, é substituída: enquanto o estudo original utiliza uma dummy binária para a existência de regras fiscais, este trabalho emprega o teto contratual de comprometimento da RCL definido na Lei n\textsuperscript{o}~9.496/1997 como variável contínua de permissividade, explorando a variação cross-sectional entre os estados. Por fim, variáveis do estudo original sem equivalente estadual direto foram omitidas, seja por indisponibilidade de dados comparáveis em nível subnacional, seja por serem, no caso de choques macroeconômicos comuns, absorvidas pelos efeitos fixos de ano.

\section{Dados e variáveis}
\label{sec:dados}

A base de dados compreende 25 unidades federativas brasileiras (24 estados e o Distrito Federal), acompanhadas em série anual de 2002 até 2023, com exclusão de Amapá e Tocantins por não terem celebrado acordos de refinanciamento sob a Lei n\textsuperscript{o}~9.496/1997.

Para o período 2002--2014, as variáveis fiscais foram recuperadas das séries históricas da Secretaria do Tesouro Nacional (STN). O endividamento (DCL e RCL) provém do Tesouro Transparente, uma vez que o SICONFI não disponibiliza o Relatório de Gestão Fiscal (RGF) antes de 2015. O resultado primário foi obtido pela base FINBRA, e os encargos da dívida, pela base de Serviço da Dívida dos Estados da STN. A partir de 2015, todas essas rubricas foram extraídas diretamente do SICONFI: DCL e RCL do RGF Anexo 2, resultado primário e encargos do Relatório Resumido da Execução Orçamentária (RREO) Anexo 6. O PIB estadual nominal provém das Contas Regionais do IBGE, disponíveis até 2023, e foi deflacionado pelo IPCA para obtenção do PIB real. Os tetos contratuais da Lei n\textsuperscript{o}~9.496/1997 foram coletados manualmente a partir dos contratos publicados no Diário Oficial da União e de relatórios do Tesouro Nacional.

A amostra inicia em 2002, primeiro ano com cobertura completa tanto para as variáveis fiscais quanto para o PIB estadual dos 25 entes federativos, período em que todos os contratos de refinanciamento da Lei n\textsuperscript{o}~9.496/1997 já estavam em plena vigência. A janela de estimação tem início em 2004 devido à defasagem de dois períodos para os instrumentos do modelo.

Como variável dependente, utiliza-se o resultado primário em proporção da Receita Corrente Líquida ($s_{it}$). A normalização pela RCL é padrão na literatura de federalismo fiscal brasileiro e assegura compatibilidade com os limites da LRF e dos contratos de refinanciamento. O endividamento, por sua vez, é mensurado pela Dívida Consolidada Líquida sobre a RCL ($d_{it}$), métrica oficial da LRF e das Resoluções do Senado Federal, que captura o endividamento líquido ao excluir ativos financeiros da dívida bruta.

A variável institucional central, o teto de comprometimento da RCL com o serviço da dívida ($teto_i$), é definida nos contratos de refinanciamento de 1997--1999 (Tabela~\ref{tab:tetos}). Nos contratos em que o limite foi estabelecido como intervalo, adotou-se o valor superior como medida de $teto_i$, privilegiando a comparabilidade entre estados. Por ser definido contratualmente antes do período amostral e não variar ao longo do tempo dentro de cada estado, o teto é tratado como característica estrutural invariante no tempo. Mudanças posteriores à Lei n\textsuperscript{o}~9.496/1997 --- como as Leis Complementares n\textsuperscript{o}~148/2014 e n\textsuperscript{o}~156/2016 --- modificaram o indexador e o prazo dos contratos, mas não alteraram as cláusulas de comprometimento máximo da RCL, que permanecem definidas em cada contrato individual celebrado ao amparo da Lei n\textsuperscript{o}~9.496/1997 durante todo o período da amostra.\footnote{O Regime de Recuperação Fiscal (LC n\textsuperscript{o}~159/2017) suspende temporariamente os pagamentos devidos à União sem alterar os limites percentuais de comprometimento da RCL fixados nos contratos originais. Aderiram ao regime dentro do período amostral o Rio de Janeiro (regime vigente desde 2017, com pagamentos suspensos por liminar), o Rio Grande do Sul (adesão aprovada em janeiro de 2022 e plano homologado em junho de 2022) e Goiás (homologação e suspensão parcial dos pagamentos em 2022). Minas Gerais teve o pedido de adesão aprovado, com plano ainda pendente de homologação ao fim da amostra. Nesses estados-anos, os pagamentos à União estiveram suspensos, ainda que o teto contratual de 1997 permanecesse formalmente definido.}

Para controle de efeitos cíclicos, utiliza-se o hiato do produto estadual ($yvar_{it}$), calculado pelo filtro Hodrick-Prescott aplicado ao logaritmo do PIB real estadual.

\section{Especificação do modelo e estratégia de estimação}
\label{sec:especificacao}

A função de reação fiscal estimada testa a condição de sustentabilidade de \citeonline{bohn1998} em contexto subnacional, condicionada ao desenho contratual da Lei n\textsuperscript{o}~9.496/1997. A equação é:
  
  \begin{equation}
s_{it} = \alpha_i + \lambda_t + \beta_1 d_{i,t-1} + \beta_2 (d_{i,t-1} \times teto_i) + \gamma \, yvar_{it} + \varepsilon_{it}
\label{eq:modelo}
\end{equation}

\noindent onde $s_{it}$ é o resultado primário sobre RCL, $d_{i,t-1}$ é a DCL/RCL defasada, $teto_i$ é o limite contratual de comprometimento da receita, $yvar_{it}$ é o hiato do produto, $\alpha_i$ são efeitos fixos de estado e $\lambda_t$ são efeitos fixos de ano.

O coeficiente $\beta_1$, combinado com $\beta_2 \times teto_i$, testa a condição de Bohn para cada estado: a reação fiscal líquida $\beta_1 + \beta_2 \times teto_i$, se positiva e significativa, indica que o resultado primário responde ao aumento do endividamento, sinalizando a presença de mecanismos de estabilização da dívida. O coeficiente $\beta_2$ é o parâmetro de interesse principal: indica se a intensidade dessa reação é condicionada pelo grau de permissividade dos contratos de refinanciamento, com significância avaliada, em cada nível de teto observado na amostra, por um teste de restrição linear.

Os efeitos fixos de estado absorvem características permanentes de cada unidade federativa, incluindo fatores que determinaram a negociação do teto contratual em 1997; os efeitos fixos de ano capturam choques macroeconômicos comuns a todos os estados (taxa de juros, inflação, câmbio, ciclo econômico nacional e mudanças institucionais federais). A interação $d_{i,t-1} \times teto_i$ sobrevive aos efeitos fixos bidirecionais porque possui variação temporal via $d_{i,t-1}$.

Embora a variação temporal de $d_{i,t-1}$ garanta a identificação da interação, a dívida defasada ainda pode ser endógena por simultaneidade dinâmica: um resultado primário fraco pode forçar mais endividamento, gerando correlação com o erro contemporâneo. Como a interação $d_{i,t-1} \times teto_i$ herda essa endogeneidade, a equação~(\ref{eq:modelo}) possui duas variáveis endógenas. Os instrumentos utilizados são $d_{i,t-2}$ e $d_{i,t-2} \times teto_i$, predeterminados em relação ao erro contemporâneo. O modelo é exatamente identificado, de modo que testes de sobreidentificação não se aplicam. A validade dos instrumentos é argumentada com base na literatura de função de reação fiscal: a dívida em $t-2$ afeta o resultado primário em $t$ predominantemente pelo canal do endividamento corrente, estratégia análoga à adotada por \citeonline{abubakarDebtReliefFiscal2025} e \citeonline{tabosa2016}. Os erros-padrão são clusterizados por estado.

A validade dessa estratégia é sustentada por três diagnósticos: a endogeneidade é verificada pelo teste de Wu-Hausman; a relevância dos instrumentos é avaliada pela estatística F do primeiro estágio, tendo como referência o critério de \citeonline{stockTestingWeakInstruments2004}; e a especificação MQO com efeitos fixos é reportada como referência para dimensionar o viés de endogeneidade.

Dois exercícios de robustez complementam a análise. O primeiro introduz uma tripla interação $d_{i,t-1} \times teto_i \times binding_{i,t-1}$, onde $binding_{i,t-1} = 1$ se a razão entre encargos\footnote{A variável de encargos foi reconstruída a partir da rubrica específica de serviço da dívida refinanciada pela Lei n\textsuperscript{o}~9.496/1997, e não do total de encargos da dívida consolidada, alinhando a medida de $binding$ ao conceito contratual de comprometimento da RCL.} e RCL no ano anterior ultrapassar o limiar prudencial de 70\% do teto contratual.\footnote{O limiar de 70\% foi adotado como critério ilustrativo para diferenciar estados operando próximos do limite de comprometimento da RCL daqueles com folga substancial, sem corresponder a um parâmetro formalmente estabelecido em norma ou contrato.} O objetivo é investigar se o efeito moderador do teto é mais forte quando o estado já operava próximo do limite. A hipótese subjacente é que a proximidade da restrição contratual amplifique o comportamento associado a tetos permissivos: estados próximos de comprometer o limite têm menor espaço de manobra e maior expectativa de renegociação futura, o que, na lógica de restrição orçamentária branda discutida na Seção~\ref{sec:regras}, reduziria ainda mais os incentivos ao ajuste preventivo --- implicando $\beta_3 < 0$. A defasagem de $binding_{i,t-1}$ preserva a predeterminação de todos os componentes da interação em relação ao erro contemporâneo, e o termo é instrumentado por $d_{i,t-2} \times teto_i \times binding_{i,t-2}$, seguindo a mesma lógica de identificação do modelo principal.

Por fim, para verificar se o resultado principal é sensível aos estados em situação de estresse fiscal agudo, a especificação principal é reestimada excluindo os estados-anos sob Regime de Recuperação Fiscal ou suspensão judicial de pagamentos --- Rio de Janeiro e Rio Grande do Sul a partir de 2017 e Goiás a partir de 2022. Esses entes concentram os episódios de maior deterioração fiscal do período e operaram, nesses anos, sob regime institucional distinto do teto contratual de 1997, o que poderia distorcer a relação estimada entre dívida e resultado primário.